% !TEX root = ../paper.tex
\section{Analysis of $\ts$ for Monotone $k$-CNFs} \label{sec:monotone}

Let $F=(X,\cC)$ be a monotone $k$-CNF where every clause has exactly $k$ literals.
Let $t=\trn(F)$ be the transversal number of $F$. We assume 
that $t\leq \frac{3n}{5}$.
Let $T$ be a transversal tree for $F$ where
$T$ is constructed using a canonical clause ordering $\Pi$.
The child edges of each of its nodes are randomly ordered from left-right independent of other nodes. Let $\pi$ denote this random tree-edge ordering. Let $r$ be its root and $Q$ its labeling.

In this section, we analyze 
$\ts$  
and prove \cref{thm:main-result} for 
the monotone case. We start
with a few ideas required to keep track of the effect of the random ordering on pruning.
We then build upon them 
\cref{sec:non-monotone} to prove \cref{thm:main-result} for general $k$-CNFs.

\subsection{Random Tree Edge Ordering and Pruning}

\begin{definition}[Cut event]
We say that an edge $e=(u,v)$ is \emph{cut} if  
$u$ has an ancestor $u'$ with a child edge $e'= (u',v')$
where $Q_{e'}=Q_e$ and 
$e'$ appears to the left of the path $P_{ru}$ according to $\pi$.
\end{definition}


We use $\cutp(e)$ to denote the event that the edge $e$ is \emph{not} cut.
We also use $\cutp(P_{uv})$ to denote the event no edge along $P_{uv}$ is cut.
We use $\cutp(u)=\cutp(P_{ru})$ to denote the event that none of the edges along the path from the root to the node $u$ are cut. 

\begin{definition}[Survival probability of paths and nodes]
The survival probability $\sigma(P_{uv})$
of a path $P_{uv}$ is $\prob(\cutp(P_{uv}))$.
The \emph{survival probability} $\sigma(u)$ of a node $u$ is $\prob(\cutp(u))$.
\end{definition}


\begin{definition}[Survival value of a transversal tree]
For a subtree $T_u$ rooted at $u$,
we define
$\total(T_u) = \sum_{\text{$v$ is leaf of $T_u$}} \sigma(v)$
as the survival value of $T_u$. We write $\total(T) = \total(T_r)$ where 
$r$ is the root node.
\end{definition}

Our main tool for upper bounding the expected running time of $\ts$ is the following basic fact.

\begin{fact}
\label{fact:ntrans-tree}
%Let $F$ be a $k$-CNF and let $T$ be a transversal tree for $F$. 
The expected number of leaves visited by $\ts$ is exactly $\total(T)$. In particular $\ntrn(F) \leq  \total(T)$.
\end{fact}

\begin{proof}
This follows by definition, since $\ts$ only visits surviving leaves of $T$ under a random tree-edge ordering. Furthermore, let $Y$ be  a minimum-size transversal of $F$.
We argue that $\sum_{v:Q_v=Y }\prob(\pi(P_{rv})) = 1$ which implies 
$\ntrn(F) \leq  \total(T)$.
\end{proof}

Our goal is to upper bound $\total(T)$ by bounding 
the survival probabilities of the leaves at depth $t$. 
A path survives if and only none of its edges are cut.
For an edge to be cut, it is necessary that some ancestor of the edge 
has a child edge with the same label as that of the edge. 
We will keep track of repeated edge labels via markings to bound 
the cut probabilities from below and thereby bounding the survival probabilities
from above.

\begin{definition}[Marking set of an  edge]
%Let $T$ be a transversal tree for $F$ with root $r$ 
%and labeling $Q$. 
The \emph{marking set} $M(e)$ of an edge $e=(u,v)$ in $T$ is the set of
nodes $w\neq u$ in the path $P_{ru}$ which have a child edge $e'$ such that $Q_e=Q_{e'}$.
We say that the nodes in $M(e)$ mark the edge $e$.
We also say that the nodes in $M(e)$ mark the label of $e$. 
%If $M(e) \neq \emptyset$, we say that the edge $e$ is marked.
\end{definition}

\begin{definition}[Marked edges]
%Let $T$ be a transversal tree. 
An edge $e=(v,u)$ in $T$ is \emph{marked}
if $M(e)\neq \emptyset$. 
\end{definition}

Marked edges are precisely those that have a non-zero probability of being cut. In fact, we can calculate the survival probability exactly.


\begin{fact}
For an edge $e$ in $T$,  $\sigma(e) = 2^{-|M(e)|}$.
\end{fact}

\begin{definition}
%Let $T$ be a transversal tree with $r$ as its root.
Let $u$ be a node in $T$. For each node $v$ 
along the path $P_{ru}$, let $N_u(v) = \{e\in P_{ru} \mid v\in M(e)\}$.
$N_u(v)$ is the set of edges along the path $P_{ru}$ marked by $v$.
\end{definition}

\begin{fact}
\label{fact:path-prob}
For any node $u$ in $T$,
the survival probability $\sigma(P_{ru})$ of the path $P_{ru}$
is given by 
\begin{align*}
\sigma(P_{ru}) = \Pi_{v: \text{ a node along }  P_{ru}} \frac{1}{|N_u(v)|+1}
\end{align*}
\end{fact}

\begin{proof}
$P_{ru}$ survives if and only if for every node $v \in P_{ru}$ and every 
edge $e \in P_{ru}$ that $v$ marks, 
the child edge of $v$ with the same label as that of $e$ appears to the right of the path. This happens with probability $\frac{1}{|N_u(v)| + 1}$ and since these events are independent, the claim follows.
\end{proof}


\subsection{An Analysis of $\ts$ for Monotone $3$-CNF}


Although \cref{fact:path-prob} gives us a fairly complete picture of 
the survival probabilities of individual paths in $T$ in terms of the edge
markings of the path,  we need a few ideas to find nontrivial upper bounds on $\total(T)$.
The first idea is the concept of a weight which captures the number of 
marked edges in a path or a shoot. 

\begin{definition}[Weight]
Let $P_{uv}$ be a path in $T$. 
The \emph{weight} of $P_{uv}$ is defined as 
$W(P_{uv}) \eqdef |\{e \in P: M(e) \ne \emptyset \}|$, 
i.e., the number of marked edges along $P_{uv}$. 
The \emph{weight} of a shoot $S_{uv}$ denoted by $W(S_{uv})$ is the number of marked edges 
in the shoot $S_{uv}$.

\end{definition}

The following fact provides a lower bound on the weight of each root to leaf
shoot  in $T$.

\begin{fact}
\label{fact:shoot-extended-weight-lower-bound}
Every root to leaf shoot $S_{ru}$ in $T$ has a weight of
at least $3t - n$.
\end{fact}


\begin{proof}
Since the depth of $T$ is $t$,
a root to leaf shoot has $3t$ edges and there are only $n$ distinct 
edge labels, at least $3t - n$ labels appear at least twice.
\end{proof}

\begin{definition}[Weight of a tree]
We say that a tree has weight $w$ if every root to leaf shoot of the tree has weight at least $w$.
\end{definition}

\begin{definition}[$M(w, d)$]
For non-negative integers 
$w$ and $d$, 
let $M(w, d)$ be the maximum survival value
over all ternary depth-$d$ transversal trees with weight $w$.
\end{definition}


We can now upper bound $\total(T)$ in terms of $M(w,d)$ 
exploiting the canonical clause ordering.
Our canonical clause ordering $\Pi$ starts with a 
maximal collection of disjoint clauses $C_1, C_2, \cdots, C_m$
so that all clauses in the formula intersect with at least one of the clauses from these disjoint clauses. 
We observe that $m\geq \frac{t}{3}$ for monotone $F$
since otherwise by setting all the variables in the monotone 
clauses $C_i$, we satisfy $F$ contradicting 
that the transversal number of $F$ is $t$.

\begin{lemma}
\label{pess-ntran}
\[
\total(T) \le
\begin{cases}
3^t & t \le \frac{n}{3}\\
3^{\frac{t}{3}}\times M(3t - n, \frac{2}{3}t) & \textrm{ otherwise}
\end{cases}
\]
\end{lemma}

\begin{proof}
%[\textbf{Proof of \cref{pess-ntran}}]
For $t\le \frac{n}{3}$, we use the trivial upper 
bound of 1 on the survival probability of paths
to conclude that $\total(T) \le 3^t$ as desired.

For $t \ge \frac{n}{3}$, we use the fact that the canonical clause
ordering starts with a maximal collection of disjoint clauses $C_1, C_2, \cdots, C_m$
where  $m\geq \frac{t}{3}$.
We observe that for $1\leq i\leq \frac{t}{3}\leq m$ 
each node at level $i$ of $T$ is expanded by the same
clause $C_i$.
Moreover, none of  the child edges  of nodes at level $1\leq i\leq \le \frac{t}{3}\leq m$ are marked
as the corresponding clauses are disjoint.
The result follows since
for every node $u$ at level $\frac{t}{3} + 1$, the subtree $T_u$ has depth $\frac{2t}{3}$ and
the weight of
every root to leaf shoot in $T_u$
is at least $3t - n$ minus the number of marked edges from root to $u$ = $3t - n - 0 = 3t - n$ 
(\cref{fact:shoot-extended-weight-lower-bound}).
\end{proof}

\subsubsection{Upper Bounds on $M(w,d)$}
%We now proceed to upper bound $M(w, d)$. 
Upper bounding $M(w,d)$ for $T$ based on random tree-edge ordering $\pi$
is challenging. 
Instead we introduce a different
random process $\pi'$ for $T$: Each edge $e$ survives with probability $p_e$ independently
where $p_e = \lambda$ if $e$ is marked and $1$ otherwise,
where we define 
$\lambda \eqdef \frac{1}{\sqrt{3}}$.
The concept of survival probability under $\pi'$
can be extended to paths and nodes of $T$.
For example, the $\sigma'(P) = \prod_{e\in P} p_e$ is the survival probability
of the path $P$ under $\pi'$. 
Similarly, we define the survival value $\total'(T)$ of the transversal tree $T$ according to $\pi'$ as $\sum_{v\text{ is a leaf}} \sigma'(P_{rv})$. 
We define $M'(w,d)$ as the maximum survival value $\total'(T')$ of transversal trees $T'$
of depth $d$ where every root to leaf shoot has weight at least $w$.
The following lemma shows that $\total(T)\leq \total'(T)$.

\begin{lemma}
\label{path-pessim}
For a root to leaf path $P_{ru}$, 
$\sigma(P) \le \lambda^{W(P_{ru})} = \sigma'(P)$ which in turn implies $\total(T)\leq \total'(T)$
and $M(w,d)\leq M'(w,d)$. 

\end{lemma}

\begin{proof}
Given an edge $e \in P_{ru}$, define the contribution $y_e$ of $e$ to the
	survival probability (according to $\pi$) of $P_{ru}$ as 
\begin{align*}
y_e \eqdef \prod_{v \in M(e)}\left(\frac{1}{|N_u(v)|+1}\right)^{1/|N_u(v)|}
\end{align*}

where the empty product is considered as 1.
By Fact \ref{fact:path-prob}, $\sigma(P_{ru}) = \prod_{e : M(e) \ne \emptyset} y_e$. It is then sufficient to show 
that $y_e \le p_e$. 
Observe that $N_u(v)$ can be at most 2 since $F$ is a $3$-CNF.
For each $v \in M(e)$ with $N_u(v) =1$, the probability that $v$ does not cut $e$
is exactly $\frac{1}{2}$, independent of other nodes.
If $N_u(v)=2$ for $v\in M(e)$, $v$ marks another edge 
$e'$ along the path in addition to $e$.
The probability that $v$ cuts neither $e$
nor $e'$ is exactly $\frac{1}{2}\times\frac{2}{3}=\frac{1}{3}$, independent of other nodes.
If $N_u(v)=2$, we regard the probability of each edge surviving 
as the geometric average $\lambda$ of the survival probabilities 
of individual edges.
As a consequence, $y_e$ can be written 
as $(\frac{1}{2})^a (\frac{1}{\sqrt{3}})^b$
for some non-negative integers $a$ and $b$ 
%such that $a + b > 0$ 
where $a$ is the number of ancestors $v$ of $e$ such 
that $v$ marks exactly one edge along the path
and $b$ is the number of ancestors $v$ of $e$ such 
that $v$ marks exactly two edges along the path.
We now argue that $y_e$ is at most $p_e$.
If $e$ is not marked, then $a+b=0$ and hence $y_e=p_e$.
If $e$ is marked, we have $a+b>0$ which implies $y_e\leq \lambda = p_e$.
\end{proof}

The following lemma determines $M'(w,d)$.

\begin{lemma} 
\label{thm:Mwd}
For all $0\le d\le n, 0\le w\le 3d$, we have
\[
M'(w,d) = 
\begin{cases}
(2+\lambda)^{w}3^{d-w} & 0\le w \le d\\
(1+2\lambda)^{w-d}(2+\lambda)^{2d-w} & d\le w\le 2d\\
(3\lambda)^{w-2d}(1+2\lambda)^{3d-w} & 2d\le w\le 3d\\
\end{cases}
\]
\end{lemma}

\cref{thm:Mwd} already gives an upper bound $\total(T) \le M'(3t - n, t)$. 
However, taking advantage 
of \cref{pess-ntran},
we can improve this bound 
to establish
\cref{thm:main-result} for monotone formulas.  

\subsubsection{Proof of \cref{thm:main-result} for monotone formulas}

\begin{proof}
%[\textbf{Proof of \cref{thm:main-result} for monotone formulas}]
\phantomsection\label{proof:monotone-main-result}
By \cref{fact:ntrans-tree} and \cref{path-pessim} the expected time of $\ts$ is bounded by 
$\total'(T)$ (up to polynomial factors). 
We divide the proof into cases based on the value of $t$.
\begin{enumerate}[label=\textbf{Case \arabic*:}, itemindent=*, leftmargin=0em]

\item
$t\le \frac{n}{3}$. Applying \cref{pess-ntran} for the case $t\le \frac{n}{3}$, we 
get $\total(T) \le 3^t$.

\item
$\frac{n}{3} < t \le \frac{3n}{7}$.
$t \le \frac{3n}{7}$ implies $3t-n \le \frac{2t}{3}$.
%Applying \cref{thm:Mwd} for the case of $w\le d$, 
%and combining it with 
We apply
\cref{pess-ntran} together with 
\cref{thm:Mwd} for the case $0\le w\le d$ to get 
\[
\total(T) \le 3^{\frac{t}{3}} M'\left(3t-n, \frac{2t}{3}\right) = \left(\frac{3}{2+\lambda}\right)^n\left(\frac{(2+\lambda)^3}{9}\right)^t \le 1.164^n \times 1.9023^t
\]

\item
$\frac{3n}{7}\le t \le \frac{n}{2}$. 
We note that $t \le \frac{n}{2}$ implies 
$3t-n \le t \le \frac{4t}{3}$ and $t \ge \frac{3n}{7}$ implies $3t-n \ge \frac{2t}{3}$. 
We apply
\cref{pess-ntran} together with 
\cref{thm:Mwd} for the case $d\le w\le 2d$ to get 
\[
\total(T) \le 3^{\frac{t}{3}} M'\left(3t-n, \frac{2t}{3}\right) = \left(\frac{2+\lambda}{1+2\lambda}\right)^n \left(\left(\frac{3(1+2\lambda)^{7}}{(2+\lambda)^{5}}\right)^{1/3}\right)^t \le 1.1962^n \times 1.7851^t
\]
\end{enumerate}
\end{proof}

\subsubsection{Proof of \cref{thm:Mwd}}

Let $T'$ be a tree of depth $d$ and weight $0\leq w\leq 3d$.
It is clear that the survival value $\total'(T')$ of $T'$ is determined once
the edges are marked consistent with the fact that every root to leaf shoot has weight at least $w$.
We say that a transversal tree $T'$ of depth $d$ and weight $w$ is \emph{normal} if it has the following marking:
Let $w=id+j$ where $0\leq i\leq 3$ and $0\leq j<d$.
Mark $i+1$ children of every non-leaf node in the 
first $j$ levels and mark $i$ children
for each of the remaining non-leaf nodes. The survival value of normal tree is exactly $((i+1)\lambda + 2 - i)^j(i\lambda + 3 - i)^{d - j}$. One can easily see that this is given exactly as $M'(w,d)$ as in the statement of the lemma.
We will show that normal trees have the largest
survival values by induction on $d$ which completes the proof
of \cref{thm:Mwd}.

Let $T'$ be a tree of depth $d$ and weight $w$ with the maximum survival value $\total'(T')$.
Let $r$ be the root of $T'$. Assume that $l_1$ children of $r$ are marked.
Let $T'_1, T'_2$, and $T'_3$ be the subtrees of $T'$ of depth $d-1$ 
and weight $w-l_1$ with $r_1, r_2$, and $r_3$ as their root nodes respectively.
$T'_1, T'_2$, and $T'_3$ are normal by induction hypothesis.
Assume that $l_2$ child edges of each of $r_i$ are marked.
The survival value of $T'$ is 
$g_1 g_2 M'(w-(l_1+l_2), d-2)$ where
$g_1= (3-l_1 +l_1 \lambda)$ and $g_2=(3-l_2 + l_2\lambda)$.
It is easy to see that if $l_1+l_2$ is held constant, $g_1 g_2$ is maximized when $l_1$ and $l_2$
are as equal as possible. If $|l_1-l_2|=1$, the survival value of $T'$ does not change if $r$ has marked $l_2$ children and 
each $r_i$ has $l_1$ marked children, that is, if the number of markings of the first two levels are exchanged.

If $l_1 = l_2$, then $T'$ is normal.
If $|l_1-l_2|\geq 2$, then $T'$ does not have the largest
survival value which is a contradiction.
We are left with the case that  $l_1$ and $l_2$ differ by one. If $l_1<l_2$, we swap $l_1$ and $l_2$ without changing the survival value and normality follows from induction. If $l_1>l_2$, the tree is already normal. Otherwise, its survival value cannot be the maximum. 
%\end{proof}


%%\begin{proof}[\textbf{Proof of \cref{thm:Mwd}}]
%We argue that the recurrence relation for $M(w,d)$ implies the bounds as claimed in
%the lemma.  The base cases can be verified easily. 
%For $d\geq1$ and $w\geq 3$, we argue as follows.
% that the recurrnce
%%implieswe guess a 
%solution for $M(w,d)$ 
%and verify it using the  recurrence in \cref{claim:Mwd-recurrence} 
%to prove \cref{thm:Mwd}.

%Let $F, F_1, F_2, F_3: \N^2\rightarrow \R$ be defined as:
%\begin{eqnarray*}
%F_1(w,d)&= & (2+\lambda)^{w}3^{d-w} \\
%F_2(w,d) & = & (1+2\lambda)^{w-d}(2+\lambda)^{2d-w} \\
%F_3(w,d)&=&   (3\lambda)^{w-2d}(1+2\lambda)^{3d-w}\\
%F(w,d) & = & \min (F_1(w,d),F_2(w,d),F_3(w,d)) .
%\end{eqnarray*}
%For $1\le i\le 3$, define $Z_i$ to be the set of pairs $(w, d)$ such that $d \ge 0$ and $w\in [(i-1)d, id]$.
%We will first argue that 
%if $(w, d)\in Z_i$, then $F(w, d) = F_i(w, d)$ 
%for $1\leq i\leq 3$.
%In other words, 
%$F_i (w,d)$ has the minimum value among the three functions
%if $(w,d) \in Z_i$ for $1\leq i\leq 3$.
%
%We show that
%$F_1(w,d) \leq F_2(w,d)$ if and only if $w \leq d$, with equality when $w=d$
%and
%$F_2(w,d) \leq F_3(w,d)$ if and only if $w \leq 2d$, with equality when $w=2d$.
%These two statements imply
%that if $(w, d)\in Z_i$, then $F(w, d) = F_i(w, d)$ 
%for $1\leq i\leq 3$.
%
%If $(w,d)\in Z_1$, it is clear 
%$\frac{F_1(w,d)}{F_2(w,d)} = 
%\left(\frac{(2+\lambda)^2}{3(1+2\lambda)}\right)^{w-d}\leq 1$
%with equality when $w=d$.
%If $(w,d)\in Z_1$, it is also the case
%$\frac{F_2(w,d)}{F_3(w,d)} = 
%\left(\frac{}{3(1+2\lambda)}\right)^{w-d}\leq 1$.
%If $(w,d)\in Z_2$, it is clear 
%$\frac{F_2(w,d)}{F_3(w,d)} = 
%\left(\frac{(1+2\lambda)^2}{3\lambda(2+\lambda)}\right)^{w-2d}\leq 1$.
%with equality when $w=2d$.
%These two statements imply
%that if $(w, d)\in Z_i$, then $F(w, d) = F_i(w, d)$ 
%for $1\leq i\leq 3$.
%%
%We will finally argue that 
%if $(w, d)\in Z_i$, then 
%$M(w,d) \le F_i(w, d)$
%for $1\leq i\leq 3$.
%
%\begin{enumerate}[label=\textbf{Case \arabic*:}, itemindent=*, leftmargin=0em]
%\item  Let $(w, d)\in Z_1$. It follows that $(w-1,d-1), (w-2,d-1), (w-3,d-1) \in Z_1$ 
%\begin{eqnarray*}
%M(w,d) & \le  & \max \{3M(w,d-1), (2+\lambda)M(w-1,d-1),(1+2\lambda)M(w-2,d-1),3\lambda M(w-3,d-1)\\
%& \leq & \max\{3F_1(w,d-1),(2+\lambda)F_1(w-1,d-1), (1+2\lambda)F_1(w-2,d-1),3\lambda F_1(w-3,d-1)\}\\
%& = & F_1(w,d) \max\{1, 1, (1+2\lambda)3/(2+\lambda)^2, 3^3\lambda/(2+\lambda)^3\}\\
%&=& F_1(w,d).
%\end{eqnarray*}
%
%\item Assume $(w, d)\in Z_2$.
%\begin{eqnarray*}
%M(w,d) & \le  & \max \{3M(w,d-1), (2+\lambda)M(w-1,d-1),(1+2\lambda)M(w-2,d-1),3\lambda M(w-3,d-1)\\
%& \leq & \max\{3F_2(w,d-1),(2+\lambda)F_2(w-1,d-1), (1+2\lambda)F_2(w-2,d-1),3\lambda F_2(w-3,d-1)\}\\
%& = & F_2(w,d) \max\{3(1+2\lambda)/(2+\lambda)^2, 1, 1, 3\lambda(2+\lambda)/(1+2\lambda)^2\}\\
%&=& F_2(w,d).
%\end{eqnarray*}
%
%\item Assume $(w, d)\in Z_3$.
%\begin{eqnarray*}
%M(w,d) & \le  & \max \{3M(w,d-1), (2+\lambda)M(w-1,d-1),(1+2\lambda)M(w-2,d-1),3\lambda M(w-3,d-1)\\
%& \leq & \max\{3F_3(w,d-1),(2+\lambda)F_3(w-1,d-1), (1+2\lambda)F_3(w-2,d-1),3\lambda F_3(w-3,d-1)\}\\
%& = & F_3(w,d) \max\{3(3\lambda)^2/(1+2\lambda)^3, (3\lambda)^2/(1+2\lambda)^2, 1, 1\}\\
%&=& F_3(w,d).
%\end{eqnarray*}
%
%\end{enumerate}
%\end{proof}
